\subsection{Canonical Penalty and Preliminary Bounds}
\label{app:canonical-prerequisites}

The fixed bounds in Section~\ref{app:fixed-arithmetic} give
\(x<1/33\), \(r>99/100\), \(\theta>3/2\), and
\(\lambda^2<P<5/2\), using \(\arcsin r<\pi r/2\).
Since \(\theta+\sinh\ell>1/(2x)\),
\begin{equation}
 \gamma>\frac{3e^\ell}{4\ell}>\ell+\frac72=2d+1.
 \label{eq:canonical-penalty-lower}
\end{equation}
The second inequality follows because \(3e^\ell-4\ell^2-14\ell\)
and its first two derivatives are positive on this half-line.

The exact identities
\begin{equation}
 \begin{split}
 p\weight(1)-b&=-\lambda^2,\\
 p\weight(2)-b&=-\frac{\lambda^2}{\gamma+2},\\
 C_d&=\lambda^2\left[d-2-\frac{(d-1)^2}{\gamma+d}\right].
 \end{split}
 \label{eq:canonical-multipliers}
\end{equation}
give \(C_d>0\), since \(\gamma(d-2)>1\), and
\(C_d<(5/2)(d-2)\), \(b<(5/2)(2+1/7)<6\).
Since \(0<1-1/(2L)<1/3\), the mean condition follows from
\begin{align*}
 \gamma-b-C_d\left(1-\frac1{2L}\right)
 &>\ell+\frac72-6-\frac56\left(\frac\ell2-\frac34\right)\\
 &=\frac16+\frac7{12}\left(\ell-\frac72\right)>0.
\end{align*}
The endpoint-principle coefficient satisfies
\begin{align*}
 \frac{p d_2}{2r}
 &=\frac p2(d-2)\frac\gamma{\gamma+d}
                   \frac\gamma{\gamma+2}r^3\\
 &>\frac3{16}>\frac5{32},
\end{align*}
using \(p/2>1\), \(d-2\ge1\), both fractions \(>1/2\), and
\(r^3>3/4\). Also \(0<d_2<d_1\), hence \(d_b>0\).

\paragraph{Entropy estimate.}
Put \(E=h(r)-L(1-r^2)\ge0\), using \(\phi(r)\le Lr^2\).
For every \(\ell\ge3/2\),
\begin{equation}
                     E\le x^2(2\ell-5/3).
 \label{eq:shared-entropy-error}
\end{equation}
Indeed, the entropy formula, \(q=x^2\), and
\(\log(1+q)/q\le(2+q)/[2(1+q)]\) give
\[
 2\ell-\frac E q-\frac53
 \ge\frac{4(L-2/3)-11q/6+5q^2/6}{(1+q)^2}>0.
\]
Its numerator exceeds \(8/81-11/120>0\), since
\(L-2/3>2/81\) and \(q<1/20\).

\subsection{Small Heights: Paired Spectral Costs}
\label{app:manual-small}

Let \(z=\ell/2\). The line-angle LSI comparison is
\begin{align*}
 M_{\rm small}&=\frac B\ell-r
       -\frac{pK_{\rm LS}r^2F(\ell)}{zF(z)},\\
 K_{\rm LS}&=\frac{\gamma^2 e^{2d-3}}{2(\gamma+d)(\gamma+1)}.
\end{align*}
Set \(a_d=d-1\), \(k=e^{-1/2}\),
\(T=e^\ell F(\ell)/F(z)\), and \(f_d=C_d/\lambda^2\).
By \eqref{eq:canonical-multipliers} and \(\phi(r)/L=r^2-E/L\),
\begin{equation}
 \begin{split}
 \frac{\ell M_{\rm small}}{\theta^2}
 ={}&d-\frac\ell{c(\ell)}-kT\\
 &-\frac{a_d(a_d-kT)}{\gamma+d}-\frac{f_dE}{Lr^2}.
 \end{split}
 \label{eq:canonical-small-pairing}
\end{equation}
The bounds \(t/(1+t)\le\log(1+t)\) and
\((1+t)\log(1+t)\le t+t^2/2\) give
\begin{align*}
 \frac{2\ell+1}{(\ell+1+x/2)(1+x)}
 &\le T\le\frac{2\ell+1+x^2/2}{(\ell+1)(1+x)}\\
 &\le2-\frac1{\ell+1}.
\end{align*}
The lower bound exceeds \(5/3\), because
\[
 \ell-2-5\ell x-\frac{15}2x-\frac52x^2
 \ge\frac32-\frac{25}{32}-\frac5{2048}>0.
\]
Here the numerator increases with \(\ell\) and decreases with \(x\).
Using \(3/5<k<5/8\) and \(c(\ell)>\theta^2>9/4\),
\begin{align*}
 1&<kT<\frac58\left(2-\frac1{\ell+1}\right),\\
 d-\frac\ell{c(\ell)}-kT
 &>\frac\ell{18}+\frac5{8(\ell+1)}\ge\frac13,
\end{align*}
the last difference from \(1/3\) being
\((2\ell-3)(2\ell-7)/[72(\ell+1)]\).
Also \(\gamma+d>4(d-1)(d-2)\): with \(s=\ell-7/2\ge0\),
the cubic Taylor lower bound gives
\begin{gather*}
 3e^\ell-(4\ell^3-6\ell^2-8\ell)\\
 >29+2s+\frac{27}2s^2+\frac{25}2s^3>0.
\end{gather*}
Together with \eqref{eq:canonical-penalty-lower}, this yields
\[
 \frac{a_d(a_d-kT)}{\gamma+d}
 <\frac{a_d(a_d-1)}{\gamma+d}<\frac14.
\]
Finally, \eqref{eq:shared-entropy-error} gives
\[
 \frac{f_dE}{Lr^2}<2(d-2)(2\ell-5/3)e^{-2\ell}<\frac1{96},
\]
using \(f_d<d-2\), \(Lr^2>1/2\), and monotonicity of the
product, whose value at \(7/2\) is below \(2(16/3)/32^2=1/96\).
Thus
\[
 M_{\rm small}>\frac{\theta^2}{\ell}
                      \left(\frac13-\frac14-\frac1{96}\right)
                =\frac{7\theta^2}{96\ell}>0.
\]
Since \(vF(v)\) decreases and \(c(v)\le v\), this covers
\(0<v\le z\) and gives \(B/\ell>r\).

\subsection{Half-Height Join and Angular Saving}
\label{app:joining-height}\label{app:manual-join-direct}

Write \(a_z=F(\ell)/F(z)\). The original Parseval margin at the
join has the exact form
\begin{align*}
 M_{\rm join}
 &:={Bz}/{\ell}-r c(z)-p[d_2/2+d_ba_z]\\
 &=r[c(\ell)-c(z)]-\frac{\lambda^2r^4}{2(\gamma+2)}
   +\frac{C_d}{2}\left[1-r^4-\frac{h(r)}L\right]\\
 &\quad-p d_ba_z.
\end{align*}
The multiplier identities also give
\begin{equation}
 \begin{split}
 p d_b={}&C_d(r^2-r^4/2)
       +\lambda^2\left[r^2-\frac{r^4}{2(\gamma+2)}\right]\\
       &\le C_d/2+P\le Pd/2.
 \end{split}
 \label{eq:canonical-db-pairing}
\end{equation}

To prove that \(e^v[P-c(v)]\) increases, write
\(t=\arcsin(\tanh v)\). Then
\begin{align*}
 &(t+\cos t)^2-c(v)-c'(v)\\
 &\quad=(1-\sin t)(t/\sin t-\cos t)^2
 +\sin t\cos^2t>0,
\end{align*}
and \(t+\cos t<\pi/2\), so \(c+c'<P\).
At \(v_0=\log(10/3)\), \(4/\pi>5/4\) and
\(\arctan(3/10)\ge3/10-(3/10)^3/3\) give
\[
 2t/\pi<16/25,\qquad
 c(v_0)/P<(16/25)^2(109/91)<1/2.
\]
Since \(z\ge7/4>v_0\), \(P-c(z)>(5P/3)\sqrt x\).
Combining this with
\(c(\ell)\ge\theta^2\ge P-2\pi x\) gives
\[
                 c(\ell)-c(z)>\frac{5P}{3}\sqrt x-2\pi x.
\]
Using \(h(r)<(2\ell+1)x^2\), \(a_z=xT<2x\),
\eqref{eq:canonical-db-pairing}, and
\eqref{eq:canonical-penalty-lower}, and dropping
\(C_d(1-r^4)/2>0\), gives
\begin{align*}
 \frac{M_{\rm join}}{P\sqrt x}
 >{}&\frac{5r}{3}
 -\left(\frac76\ell+\frac54+\frac8\pi\right)e^{-\ell/2}\\
 &-\frac{(d-2)(2\ell+1)}{2L}e^{-3\ell/2}.
\end{align*}
The first term increases and both losses decrease. At \(7/2\),
\(r>99/100\), \(\sqrt x<7/40\), \(8/\pi<8/3\), and
\(2L>4/3\) bound the losses by \(8(7/40)\) and
\(6(7/40)^3<1/25\). Hence
\[
             \frac{M_{\rm join}}{P\sqrt x}
                  >\frac{33}{20}-\frac75-\frac1{25}
                   =\frac{21}{100}>0.
\]

\subsection{Cutoff Bound with Exact Low-Degree Multipliers}
\label{app:manual-cutoff-common}\label{app:pivotal-tail-extension}

Let \(F(V)=F(\ell)/a_{\rm cut}\). By
\eqref{eq:analytic-profile-speed},
\[
                  0<\ell-V<\frac53x.
\]
The Parseval margin is
\[
 M_{\rm cut}=\frac B\ell V-r c(V)
                  -p[d_2/2+d_ba_{\rm cut}].
\]
The identity
\[
 B-pd_1-r c(\ell)=C_d[\phi(r)/L-r^2]=-C_dE/L
\]
and \(c'>0\) give
\[
 M_{\rm cut}\ge-C_dE/L+p d_b(1-a_{\rm cut})
                              -\frac B\ell(\ell-V).
\]
Using \(B\le2\lambda^2r^2+C_d\), the equality in
\eqref{eq:canonical-db-pairing}, and \eqref{eq:shared-entropy-error},
\begin{align*}
 \frac{M_{\rm cut}}x>{}&C_d\left[
 \frac{8r^2(1-r^2/2)}{3+8x}-\frac{(2\ell-5/3)x}{L}
                                      -\frac5{3\ell}\right]\\
 &+\lambda^2r^2\left[
 \frac8{3+8x}\left(1-\frac{r^2}{2(\gamma+2)}\right)
                                      -\frac{10}{3\ell}\right].
\end{align*}
Since \(r^2>3/4\), \(x<1/32\), and
\((2\ell-5/3)e^{-\ell}\) decreases,
\begin{align*}
 \frac{8r^2(1-r^2/2)}{3+8x}&>\frac{32}{13}\frac{15}{32}>1,\\
 \frac{(2\ell-5/3)x}{L}&<\frac{16}{3}\frac1{32}\frac32=\frac14,\\
 \frac5{3\ell}&\le\frac{10}{21}<\frac12.
\end{align*}
Thus the first bracket exceeds \(1/4\); the second exceeds
\(2(3/4)-1=1/2\), using \(\gamma>7\) and
\(10/(3\ell)\le20/21<1\). Therefore
\[
              M_{\rm cut}>x\left(\frac{C_d}{4}
                                  +\frac{\lambda^2r^2}{2}\right)>0.
\]
Together with the local cutoff and the endpoint principle, this
completes the height comparison for every \(\ell\ge7/2\).

